% =====================================================================
% TEMPLATE — Conclusions report (methodology asset A-P10-tpl)
% =====================================================================
% Purpose: the phase-P10 deliverable — a self-contained report that documents
% the methods built for the problem, presents the P9 experimental comparison,
% draws honest conclusions, and closes with a methodology retrospective.
%
% Fill every <PLACEHOLDER>; delete guidance comments. Meant to be \input from a
% main.tex (typically LAST, after the formulation part), whose preamble provides
% at least:
%     \usepackage{amsmath,amssymb}
% (see the working harness in docs/<problem>/ — main.tex + compile.py).
%
% Structure (keep the skeleton stable; enrich guidance only):
%   Overview → Methods (exact + heuristic) → Experimental setup →
%   Results (table) → Discussion (with honest caveats) →
%   Conclusions & recommendations → Methodology retrospective.
%
% Principles this report must honor:
%   * Report ONLY what the P9 battery actually produced — no unrun claims.
%   * Every reported solution passed the P3 checker (compliance gate); say so.
%   * Use ONE fixed gap convention, defined once and reused (match P9).
%   * Be honest about the benchmark's limits (size, coverage, unclosed runs).
% =====================================================================

% =====================================================================
\part{Report --- methods, results, conclusions}
% =====================================================================

\section{Overview}
% 3–5 sentences: restate the problem and its single headline objective, name the
% techniques compared (typically one EXACT method + one HEURISTIC), and state the
% shared contract that makes the comparison fair (same solution schema, every
% solution validated by the P3 checker before it counts). No results yet.
<OVERVIEW: what problem is solved, the headline objective, the <N> techniques
compared, and the shared solution-schema + checker contract>

\section{Methods}

\subsection{Exact --- <EXACT METHOD NAME>}
% One paragraph. Summarize the exact model (reference the formulation part rather
% than repeating it): the core decision variables/objects, the objective, the
% governing constraint families, and the solver stack + the run knobs used
% (time limit, gap tolerance). Keep it to essentials — the full model lives in
% the formulation part.
<EXACT-METHOD SUMMARY: model in brief, objective, key constraints, solver stack,
run knobs (time limit / gap)>

\subsection{Heuristic --- <HEURISTIC METHOD NAME>}
% Name the approach and, if applicable, cite the phase-P7 candidate it realizes.
% Describe its components as a short numbered list. A common shape:
\begin{enumerate}
  \item \textbf{<CONSTRUCTIVE STEP>} --- <the requirement-respecting greedy that
        builds feasible building blocks; note which hard requirements it enforces
        BY CONSTRUCTION, referencing their [RQxx] ids>.
  \item \textbf{<OBJECTIVE EVALUATOR>} --- <the deterministic evaluator that turns
        a candidate into a FULL, checker-valid solution and scores it; note this
        is the reusable core the experiment harness measures>.
  \item \textbf{<SEARCH / METAHEURISTIC>} --- <the improvement loop: the degree of
        freedom perturbed, the accept rule, restarts>.
\end{enumerate}

% One closing sentence: both techniques return solutions in the same P3 schema,
% so the checker and the experiment harness treat them identically.
<CLOSING: both methods emit the same solution schema, judged identically>

\section{Experimental setup}
% Describe the benchmark, harness, compliance gate, gap convention and
% reproducibility — enough for a reader to trust and reproduce the numbers.

% Benchmark: a deterministic grid of instances from the P2 generator (seed-first),
% sized so the exact method mostly provides a trustworthy baseline. Present the
% grid as a compact table of the size dimensions that matter for this problem.
The benchmark is a deterministic grid of <N> instances from the P2 generator
(seed-first), sized so the exact method mostly <closes / proves optimality>
(the reference baseline):

\begin{center}
\begin{tabular}{l<COLUMN SPEC: one r per size dimension>}
\hline
instance & <SIZE DIM 1> & <SIZE DIM 2> & <...> \\
\hline
<INSTANCE FAMILY 1> & <val> & <val> & <...> \\
<INSTANCE FAMILY 2> & <val> & <val> & <...> \\
% ... one row per instance family (group seeds compactly, e.g. seed\{1,2,3\})
\hline
\end{tabular}
\end{center}

% Harness + gates + reproducibility. State: dual logging (rich per-run JSON +
% append-only CSV ledger), the compliance gate (every solution through the P3
% checker; non-compliant never counts), reproducibility stamps (seed, config,
% git commit, host, solver/library versions), and the ONE gap convention (define
% it explicitly, with the objective direction that makes a positive value mean
% "method beats the baseline").
Every run is dual-logged (rich per-run JSON + an append-only CSV ledger) with
reproducibility stamps (seed, config, git commit, host, solver versions). Each
solution passes the P3 checker (compliance gate). The comparison uses the fixed
gap $=(\text{baseline}_{\text{obj}}-\text{method}_{\text{obj}})/\text{baseline}_{\text{obj}}$
% ^ flip the numerator sign if the objective is MAXIMIZED; state what a positive
%   value means for THIS objective direction.
(<what a positive value means for this objective's direction>).

\section{Results}
% One row per instance. Columns: the baseline (exact) objective, its status
% (optimal / hit-time-limit / ...) and wall time, then each method's objective and
% wall time. Keep numbers exactly as logged in the CSV ledger.

\begin{center}
\begin{tabular}{l<COLUMN SPEC>}
\hline
instance & <BASELINE> obj & <BASELINE> status & <BASELINE>\,s & <METHOD> obj & <METHOD>\,s \\
\hline
<INSTANCE> & <obj> & <status> & <time> & <obj> & <time> \\
% ... one row per instance
\hline
\end{tabular}
\end{center}

% Narrative under the table: the headline finding stated plainly. Report the gap
% across the benchmark, how many instances the exact method closed, and any
% instance where it stalled and how the heuristic compared there (objective +
% speed factor). Reaffirm that every reported solution is checker-valid and give
% the heuristic's runtime envelope.
<RESULTS NARRATIVE: headline gap across the benchmark; exact-method closed count;
behavior on any unclosed instance (objective tie/beat + speed factor); all
solutions checker-valid; runtime envelope>

\section{Discussion}
% Interpret the results, then state the HONEST CAVEATS explicitly — this section
% earns the report its credibility. Typical caveats to consider (keep those that
% apply, add problem-specific ones):
%   (i)  benchmark size/coverage — is this a correctness validation or a scaling
%        study? say which it is NOT yet.
%   (ii) did the exact method actually get stressed? if few/no instances forced a
%        strictly positive gap, say the margin was not demonstrated (not that the
%        method cannot win).
%  (iii) any determinism/variance limits (single seed per method, RNG, etc.).
<DISCUSSION: interpretation of the comparison, followed by the honest caveats
that bound what these results do and do not establish>

\section{Conclusions and recommendations}
% 2–4 sentences. State which method(s) to use and why, grounded strictly in the
% evidence above. Then give the concrete recommended NEXT step(s) (e.g. a scaling
% study on larger instances where the exact method no longer baselines,
% parameter sweeps, multi-seed variance).
<CONCLUSIONS: recommended method(s) and why, plus the concrete next step(s)>

\section{Methodology retrospective}
% Step out of the domain and reflect on the METHOD itself (this is the P10
% synthesis half). Two short lists:
%   * What worked: which phases / reusable assets paid off (e.g. the phase
%     pipeline, the P3 checker reused as a compliance gate across phases, the
%     systematic experiment harness, the P7 method review that surfaced the
%     chosen heuristic).
%   * To improve: concrete, actionable gaps observed while running the case
%     (benchmark growth, storing instances as files, template generalizations a
%     second problem would need, etc.). Feed these back as decision-log candidates.
<RETROSPECTIVE — What worked: <phases/assets that paid off>.
To improve: <concrete methodological gaps observed and how to address them>>
